\chapter{Unit 6: Embedded System and Its Application}

\section{Section A: Multiple Choice Question Practice (60 Curated Problems)}

\noindent\textbf{Q1. What is an embedded system?} \hfill [Topic: introduction to embedded system | Easy]
\begin{enumerate}[label=(\alph*),topsep=2pt,itemsep=1pt,partopsep=0pt,parsep=0pt]
    \item {A type of operating system for servers}
    \item {A general-purpose computer used for office work}
    \item {A dedicated computer system designed to perform a specific task}
    \item {A software used only for gaming}
\end{enumerate}
\vspace{2mm}

\noindent\textbf{Q2. Which of the following is a common example of an embedded system?} \hfill [Topic: introduction to embedded system | Easy]
\begin{enumerate}[label=(\alph*),topsep=2pt,itemsep=1pt,partopsep=0pt,parsep=0pt]
    \item {A spreadsheet application}
    \item {A desktop PC}
    \item {A cloud data center}
    \item {A washing machine controller}
\end{enumerate}
\vspace{2mm}

\noindent\textbf{Q3. Which component acts as the 'brain' of an embedded system?} \hfill [Topic: introduction to embedded system | Easy]
\begin{enumerate}[label=(\alph*),topsep=2pt,itemsep=1pt,partopsep=0pt,parsep=0pt]
    \item {Plastic casing}
    \item {Processor (microcontroller/microprocessor)}
    \item {Display screen}
    \item {Power cable}
\end{enumerate}
\vspace{2mm}

\noindent\textbf{Q4. Embedded systems are generally designed to be:} \hfill [Topic: introduction to embedded system | Easy]
\begin{enumerate}[label=(\alph*),topsep=2pt,itemsep=1pt,partopsep=0pt,parsep=0pt]
    \item {Compact, reliable, and low power}
    \item {Large and expensive}
    \item {Only used in laboratories}
    \item {Incapable of real-time operation}
\end{enumerate}
\vspace{2mm}

\noindent\textbf{Q5. Which device integrates CPU, memory, and I/O ports on a single chip?} \hfill [Topic: comparison of microprocessor and micro-controller | Easy]
\begin{enumerate}[label=(\alph*),topsep=2pt,itemsep=1pt,partopsep=0pt,parsep=0pt]
    \item {Capacitor}
    \item {Microcontroller}
    \item {Microprocessor}
    \item {Transistor}
\end{enumerate}
\vspace{2mm}

\noindent\textbf{Q6. A microprocessor typically requires which of the following to function?} \hfill [Topic: comparison of microprocessor and micro-controller | Easy]
\begin{enumerate}[label=(\alph*),topsep=2pt,itemsep=1pt,partopsep=0pt,parsep=0pt]
    \item {External memory and peripherals}
    \item {Only a display}
    \item {No additional components at all}
    \item {Only a battery}
\end{enumerate}
\vspace{2mm}

\noindent\textbf{Q7. Which is generally more suitable for small, dedicated control applications?} \hfill [Topic: comparison of microprocessor and micro-controller | Easy]
\begin{enumerate}[label=(\alph*),topsep=2pt,itemsep=1pt,partopsep=0pt,parsep=0pt]
    \item {Supercomputer}
    \item {Microprocessor}
    \item {Mainframe}
    \item {Microcontroller}
\end{enumerate}
\vspace{2mm}

\noindent\textbf{Q8. Compared to a microcontroller, a microprocessor is usually used in:} \hfill [Topic: comparison of microprocessor and micro-controller | Easy]
\begin{enumerate}[label=(\alph*),topsep=2pt,itemsep=1pt,partopsep=0pt,parsep=0pt]
    \item {Only in sensors}
    \item {Simple toys only}
    \item {High-performance computing systems like PCs}
    \item {Only in LED lights}
\end{enumerate}
\vspace{2mm}

\noindent\textbf{Q9. Which generally consumes less power for embedded tasks?} \hfill [Topic: comparison of microprocessor and micro-controller | Easy]
\begin{enumerate}[label=(\alph*),topsep=2pt,itemsep=1pt,partopsep=0pt,parsep=0pt]
    \item {Both consume equal very high power}
    \item {Microprocessor}
    \item {Neither consumes any power}
    \item {Microcontroller}
\end{enumerate}
\vspace{2mm}

\noindent\textbf{Q10. Arduino is best described as a/an:} \hfill [Topic: introduction of arduino and sensors | Easy]
\begin{enumerate}[label=(\alph*),topsep=2pt,itemsep=1pt,partopsep=0pt,parsep=0pt]
    \item {Cloud storage service}
    \item {High-voltage transformer}
    \item {Open-source microcontroller-based development board}
    \item {Operating system for laptops}
\end{enumerate}
\vspace{2mm}

\noindent\textbf{Q11. What is the main function of a sensor?} \hfill [Topic: introduction of arduino and sensors | Easy]
\begin{enumerate}[label=(\alph*),topsep=2pt,itemsep=1pt,partopsep=0pt,parsep=0pt]
    \item {To display video output}
    \item {To supply high voltage}
    \item {To convert a physical quantity into an electrical signal}
    \item {To store large amounts of data permanently}
\end{enumerate}
\vspace{2mm}

\noindent\textbf{Q12. Which of the following is an example of a temperature sensor?} \hfill [Topic: introduction of arduino and sensors | Easy]
\begin{enumerate}[label=(\alph*),topsep=2pt,itemsep=1pt,partopsep=0pt,parsep=0pt]
    \item {LM35}
    \item {LED}
    \item {Resistor}
    \item {Relay}
\end{enumerate}
\vspace{2mm}

\noindent\textbf{Q13. The programming environment used to write code for Arduino is called the:} \hfill [Topic: introduction of arduino and sensors | Easy]
\begin{enumerate}[label=(\alph*),topsep=2pt,itemsep=1pt,partopsep=0pt,parsep=0pt]
    \item {Excel Macro}
    \item {Photoshop}
    \item {Arduino IDE}
    \item {Microsoft Word}
\end{enumerate}
\vspace{2mm}

\noindent\textbf{Q14. Arduino boards are commonly programmed using a language based on:} \hfill [Topic: introduction of arduino and sensors | Easy]
\begin{enumerate}[label=(\alph*),topsep=2pt,itemsep=1pt,partopsep=0pt,parsep=0pt]
    \item {Assembly only}
    \item {SQL}
    \item {C/C++}
    \item {HTML}
\end{enumerate}
\vspace{2mm}

\noindent\textbf{Q15. What does IoT stand for?} \hfill [Topic: role of IOT and cloud computing in condition monitoring of plant processes, health care, agriculture, manufacturing, automobiles and smart grid | Easy]
\begin{enumerate}[label=(\alph*),topsep=2pt,itemsep=1pt,partopsep=0pt,parsep=0pt]
    \item {Input of Technology}
    \item {Internet of Transistors}
    \item {Integration of Tools}
    \item {Internet of Things}
\end{enumerate}
\vspace{2mm}

\noindent\textbf{Q16. In healthcare, IoT devices are commonly used to:} \hfill [Topic: role of IOT and cloud computing in condition monitoring of plant processes, health care, agriculture, manufacturing, automobiles and smart grid | Easy]
\begin{enumerate}[label=(\alph*),topsep=2pt,itemsep=1pt,partopsep=0pt,parsep=0pt]
    \item {Remotely monitor patient vital signs}
    \item {Replace all doctors permanently}
    \item {Generate electricity for hospitals}
    \item {Store printed paper records}
\end{enumerate}
\vspace{2mm}

\noindent\textbf{Q17. In agriculture, IoT sensors are mainly used to monitor:} \hfill [Topic: role of IOT and cloud computing in condition monitoring of plant processes, health care, agriculture, manufacturing, automobiles and smart grid | Easy]
\begin{enumerate}[label=(\alph*),topsep=2pt,itemsep=1pt,partopsep=0pt,parsep=0pt]
    \item {Internet browsing history}
    \item {Soil moisture and crop conditions}
    \item {Television channels}
    \item {Mobile phone games}
\end{enumerate}
\vspace{2mm}

\noindent\textbf{Q18. What is the main role of cloud computing in IoT systems?} \hfill [Topic: role of IOT and cloud computing in condition monitoring of plant processes, health care, agriculture, manufacturing, automobiles and smart grid | Easy]
\begin{enumerate}[label=(\alph*),topsep=2pt,itemsep=1pt,partopsep=0pt,parsep=0pt]
    \item {Storing and processing data from connected devices}
    \item {Manufacturing microcontrollers}
    \item {Supplying electrical power to devices}
    \item {Generating physical sensors}
\end{enumerate}
\vspace{2mm}

\noindent\textbf{Q19. A smart grid primarily uses IoT and sensors to:} \hfill [Topic: role of IOT and cloud computing in condition monitoring of plant processes, health care, agriculture, manufacturing, automobiles and smart grid | Easy]
\begin{enumerate}[label=(\alph*),topsep=2pt,itemsep=1pt,partopsep=0pt,parsep=0pt]
    \item {Monitor and manage electricity distribution efficiently}
    \item {Stop electricity usage completely}
    \item {Produce coal for power plants}
    \item {Replace all power lines with batteries}
\end{enumerate}
\vspace{2mm}

\noindent\textbf{Q20. Condition monitoring of plant machinery using IoT helps mainly to:} \hfill [Topic: role of IOT and cloud computing in condition monitoring of plant processes, health care, agriculture, manufacturing, automobiles and smart grid | Easy]
\begin{enumerate}[label=(\alph*),topsep=2pt,itemsep=1pt,partopsep=0pt,parsep=0pt]
    \item {Detect faults early and reduce downtime}
    \item {Remove the need for electricity}
    \item {Make machines heavier}
    \item {Increase the size of machines}
\end{enumerate}
\vspace{2mm}

\noindent\textbf{Q21. A digital washing machine controller is dedicated to running only wash cycles and cannot be reprogrammed by the user for other tasks. Which characteristic of embedded systems does this best illustrate?} \hfill [Topic: introduction to embedded system | Medium]
\begin{enumerate}[label=(\alph*),topsep=2pt,itemsep=1pt,partopsep=0pt,parsep=0pt]
    \item {Application-specific and dedicated functionality}
    \item {Unlimited reconfigurable software}
    \item {General-purpose computing capability}
    \item {High-latency batch processing}
\end{enumerate}
\vspace{2mm}

\noindent\textbf{Q22. Which combination of constraints is most typical when designing a battery-powered embedded device such as a fitness band?} \hfill [Topic: introduction to embedded system | Medium]
\begin{enumerate}[label=(\alph*),topsep=2pt,itemsep=1pt,partopsep=0pt,parsep=0pt]
    \item {No timing constraints and large size}
    \item {Low power consumption and real-time response}
    \item {Maximum cost and manual operation}
    \item {High power draw and unlimited memory}
\end{enumerate}
\vspace{2mm}

\noindent\textbf{Q23. A real-time embedded system in an airbag controller must respond within a strict deadline. Missing this deadline causes system failure. This is an example of which type of system?} \hfill [Topic: introduction to embedded system | Medium]
\begin{enumerate}[label=(\alph*),topsep=2pt,itemsep=1pt,partopsep=0pt,parsep=0pt]
    \item {Batch processing system}
    \item {Hard real-time system}
    \item {Soft real-time system}
    \item {Non-real-time system}
\end{enumerate}
\vspace{2mm}

\noindent\textbf{Q24. An engineer needs a single chip containing CPU, RAM, ROM, and I/O ports for a low-cost temperature logger. Which device is the most suitable choice?} \hfill [Topic: comparison of microprocessor and micro-controller | Medium]
\begin{enumerate}[label=(\alph*),topsep=2pt,itemsep=1pt,partopsep=0pt,parsep=0pt]
    \item {Microcontroller}
    \item {Discrete logic gates}
    \item {Microprocessor}
    \item {External DRAM module}
\end{enumerate}
\vspace{2mm}

\noindent\textbf{Q25. Why does a microprocessor-based system generally require more external components than a microcontroller-based system?} \hfill [Topic: comparison of microprocessor and micro-controller | Medium]
\begin{enumerate}[label=(\alph*),topsep=2pt,itemsep=1pt,partopsep=0pt,parsep=0pt]
    \item {The microprocessor has slower clock speeds}
    \item {The microprocessor consumes no power}
    \item {Memory and I/O are not integrated on the microprocessor chip}
    \item {The microprocessor cannot execute arithmetic}
\end{enumerate}
\vspace{2mm}

\noindent\textbf{Q26. Which statement best compares typical use cases of microprocessors and microcontrollers?} \hfill [Topic: comparison of microprocessor and micro-controller | Medium]
\begin{enumerate}[label=(\alph*),topsep=2pt,itemsep=1pt,partopsep=0pt,parsep=0pt]
    \item {Microprocessors suit dedicated control; microcontrollers suit general computing}
    \item {Both are unable to interface with sensors}
    \item {Both are used only in desktop computers}
    \item {Microprocessors suit computation-heavy general systems; microcontrollers suit dedicated control tasks}
\end{enumerate}
\vspace{2mm}

\noindent\textbf{Q27. A microcontroller is often said to follow a Harvard architecture while many microprocessors use von Neumann architecture. What is the key implication of Harvard architecture?} \hfill [Topic: comparison of microprocessor and micro-controller | Medium]
\begin{enumerate}[label=(\alph*),topsep=2pt,itemsep=1pt,partopsep=0pt,parsep=0pt]
    \item {Elimination of the CPU}
    \item {Separate buses for program and data memory}
    \item {No ability to store programs}
    \item {A single shared bus for all memory}
\end{enumerate}
\vspace{2mm}

\noindent\textbf{Q28. On an Arduino Uno, which function is used to read a varying analog voltage from a sensor such as an LDR?} \hfill [Topic: introduction of arduino and sensors | Medium]
\begin{enumerate}[label=(\alph*),topsep=2pt,itemsep=1pt,partopsep=0pt,parsep=0pt]
    \item {pinMode()}
    \item {analogRead()}
    \item {digitalWrite()}
    \item {Serial.print()}
\end{enumerate}
\vspace{2mm}

\noindent\textbf{Q29. The Arduino Uno's analog-to-digital converter has 10-bit resolution. How many discrete levels can it represent?} \hfill [Topic: introduction of arduino and sensors | Medium]
\begin{enumerate}[label=(\alph*),topsep=2pt,itemsep=1pt,partopsep=0pt,parsep=0pt]
    \item {512}
    \item {4096}
    \item {1024}
    \item {256}
\end{enumerate}
\vspace{2mm}

\noindent\textbf{Q30. A designer wants to measure distance to an obstacle for an automatic parking system. Which sensor is most appropriate?} \hfill [Topic: introduction of arduino and sensors | Medium]
\begin{enumerate}[label=(\alph*),topsep=2pt,itemsep=1pt,partopsep=0pt,parsep=0pt]
    \item {LDR (light sensor)}
    \item {Thermistor}
    \item {Hall-effect switch}
    \item {Ultrasonic sensor}
\end{enumerate}
\vspace{2mm}

\noindent\textbf{Q31. Which classification best describes a thermocouple used with an Arduino to measure temperature?} \hfill [Topic: introduction of arduino and sensors | Medium]
\begin{enumerate}[label=(\alph*),topsep=2pt,itemsep=1pt,partopsep=0pt,parsep=0pt]
    \item {Analog sensor producing a continuous voltage}
    \item {Actuator converting electrical to motion}
    \item {Communication module}
    \item {Digital sensor producing only 0 or 1}
\end{enumerate}
\vspace{2mm}

\noindent\textbf{Q32. In an Arduino sketch, which function runs continuously and repeatedly after startup to handle ongoing sensor polling?} \hfill [Topic: introduction of arduino and sensors | Medium]
\begin{enumerate}[label=(\alph*),topsep=2pt,itemsep=1pt,partopsep=0pt,parsep=0pt]
    \item {setup()}
    \item {loop()}
    \item {main()}
    \item {init()}
\end{enumerate}
\vspace{2mm}

\noindent\textbf{Q33. In an IoT-based plant condition-monitoring system, vibration sensors on a motor send data to the cloud, which detects a rising trend before failure. This capability is best described as:} \hfill [Topic: role of IOT and cloud computing in condition monitoring of plant processes | Medium]
\begin{enumerate}[label=(\alph*),topsep=2pt,itemsep=1pt,partopsep=0pt,parsep=0pt]
    \item {Random fault guessing}
    \item {Reactive repair after breakdown}
    \item {Predictive maintenance}
    \item {Manual periodic inspection only}
\end{enumerate}
\vspace{2mm}

\noindent\textbf{Q34. Why is cloud computing advantageous for storing and analyzing data from thousands of plant sensors?} \hfill [Topic: role of IOT and cloud computing in condition monitoring of plant processes | Medium]
\begin{enumerate}[label=(\alph*),topsep=2pt,itemsep=1pt,partopsep=0pt,parsep=0pt]
    \item {It offers scalable storage and processing on demand}
    \item {It permanently stops all machine faults}
    \item {It eliminates the need for any sensors}
    \item {It requires no network connectivity}
\end{enumerate}
\vspace{2mm}

\noindent\textbf{Q35. An IoT monitoring system needs instant local reaction to a dangerous pressure spike without waiting for the cloud. Which approach addresses this requirement?} \hfill [Topic: role of IOT and cloud computing in condition monitoring of plant processes | Medium]
\begin{enumerate}[label=(\alph*),topsep=2pt,itemsep=1pt,partopsep=0pt,parsep=0pt]
    \item {Sending all data only to a distant cloud}
    \item {Monthly manual data upload}
    \item {Disabling the sensors}
    \item {Edge computing near the sensor}
\end{enumerate}
\vspace{2mm}

\noindent\textbf{Q36. A wearable IoT device continuously transmits a patient's heart rate to doctors for remote observation. This application is best known as:} \hfill [Topic: health care | Medium]
\begin{enumerate}[label=(\alph*),topsep=2pt,itemsep=1pt,partopsep=0pt,parsep=0pt]
    \item {Remote patient monitoring}
    \item {In-person surgery}
    \item {Hospital catering}
    \item {Pharmacy billing}
\end{enumerate}
\vspace{2mm}

\noindent\textbf{Q37. A smart-farming IoT system uses soil-moisture sensors to turn on irrigation only when moisture drops below a threshold. What is the primary benefit?} \hfill [Topic: agriculture | Medium]
\begin{enumerate}[label=(\alph*),topsep=2pt,itemsep=1pt,partopsep=0pt,parsep=0pt]
    \item {Increased water wastage}
    \item {Removal of all crops}
    \item {Efficient water usage}
    \item {Higher fuel consumption}
\end{enumerate}
\vspace{2mm}

\noindent\textbf{Q38. In a smart manufacturing line, IoT sensors feed real-time production data to a dashboard that automatically flags defective units. This is a core concept of:} \hfill [Topic: manufacturing | Medium]
\begin{enumerate}[label=(\alph*),topsep=2pt,itemsep=1pt,partopsep=0pt,parsep=0pt]
    \item {Industry 4.0}
    \item {Manual assembly only}
    \item {Paper-based record keeping}
    \item {Pre-electricity production}
\end{enumerate}
\vspace{2mm}

\noindent\textbf{Q39. A connected car uploads engine diagnostics to the cloud, which alerts the driver about an impending component failure. This feature primarily relies on:} \hfill [Topic: automobiles | Medium]
\begin{enumerate}[label=(\alph*),topsep=2pt,itemsep=1pt,partopsep=0pt,parsep=0pt]
    \item {IoT telematics and cloud analytics}
    \item {A standalone offline odometer}
    \item {Printed service manuals}
    \item {Mechanical gear ratios only}
\end{enumerate}
\vspace{2mm}

\noindent\textbf{Q40. A smart grid uses IoT-enabled smart meters that allow two-way communication between consumers and the utility. What is a key advantage of this over traditional meters?} \hfill [Topic: smart grid | Medium]
\begin{enumerate}[label=(\alph*),topsep=2pt,itemsep=1pt,partopsep=0pt,parsep=0pt]
    \item {Reading only once per year manually}
    \item {Inability to measure consumption}
    \item {No communication with the utility}
    \item {Real-time demand response and dynamic load management}
\end{enumerate}
\vspace{2mm}

\noindent\textbf{Q41. A real-time embedded system must guarantee a response to a sensor interrupt within 50 $\mu$ s . The CPU runs at 20 MHz and the interrupt latency plus the ISR consumes 600 clock cycles. By what minimum factor must the clock be increased to safely meet a stricter deadline of 20 $\mu$ s ?} \hfill [Topic: introduction to embedded system | Hard]
\begin{enumerate}[label=(\alph*),topsep=2pt,itemsep=1pt,partopsep=0pt,parsep=0pt]
    \item {1.5 $\times$}
    \item {2.5 $\times$}
    \item {3 $\times$}
    \item {5 $\times$}
\end{enumerate}
\vspace{2mm}

\noindent\textbf{Q42. Which characteristic most sharply distinguishes a hard real-time embedded system from a soft real-time embedded system?} \hfill [Topic: introduction to embedded system | Hard]
\begin{enumerate}[label=(\alph*),topsep=2pt,itemsep=1pt,partopsep=0pt,parsep=0pt]
    \item {It always uses a faster processor than soft real-time systems}
    \item {It never uses interrupts for timing control}
    \item {Missing a deadline causes total system failure, not just degraded quality}
    \item {It stores more data in non-volatile memory}
\end{enumerate}
\vspace{2mm}

\noindent\textbf{Q43. A design requires on-chip RAM, ROM, timers, and ADC in a single IC for a low-cost washing machine controller. Choosing a microprocessor instead of a microcontroller would MOST directly cause which penalty?} \hfill [Topic: comparison of microprocessor and micro-controller | Hard]
\begin{enumerate}[label=(\alph*),topsep=2pt,itemsep=1pt,partopsep=0pt,parsep=0pt]
    \item {Increased board area, cost, and power due to external peripheral chips}
    \item {Loss of the ability to run assembly language}
    \item {Reduced instruction set width limiting arithmetic}
    \item {Inability to execute any interrupt service routine}
\end{enumerate}
\vspace{2mm}

\noindent\textbf{Q44. The 8051 microcontroller uses a Harvard architecture while a typical general-purpose microprocessor uses von Neumann. What is the key performance consequence of this difference?} \hfill [Topic: comparison of microprocessor and micro-controller | Hard]
\begin{enumerate}[label=(\alph*),topsep=2pt,itemsep=1pt,partopsep=0pt,parsep=0pt]
    \item {Simultaneous instruction fetch and data access improve throughput}
    \item {Interrupts are completely eliminated}
    \item {Program memory can never be larger than data memory}
    \item {Only one addressing mode becomes available}
\end{enumerate}
\vspace{2mm}

\noindent\textbf{Q45. A system must perform heavy floating-point image processing with a large external DRAM and run a full OS. Which selection and reasoning is correct?} \hfill [Topic: comparison of microprocessor and micro-controller | Hard]
\begin{enumerate}[label=(\alph*),topsep=2pt,itemsep=1pt,partopsep=0pt,parsep=0pt]
    \item {Microcontroller, because its on-chip RAM handles large images}
    \item {Microprocessor, because it scales with large external memory and OS support}
    \item {Microprocessor, because it cannot interface external memory}
    \item {Microcontroller, because it has a built-in FPU superior to any microprocessor}
\end{enumerate}
\vspace{2mm}

\noindent\textbf{Q46. An Arduino Uno (ATmega328P) uses a 10 -bit ADC with a 5 V reference. A sensor outputs 2.5 V . What raw ADC value and resolution per step apply?} \hfill [Topic: introduction of arduino and sensors | Hard]
\begin{enumerate}[label=(\alph*),topsep=2pt,itemsep=1pt,partopsep=0pt,parsep=0pt]
    \item {Value $\approx$ 512 , step $\approx$ 9.77 mV}
    \item {Value $\approx$ 512 , step $\approx$ 4.88 mV}
    \item {Value $\approx$ 1023 , step $\approx$ 4.88 mV}
    \item {Value $\approx$ 256 , step $\approx$ 4.88 mV}
\end{enumerate}
\vspace{2mm}

\noindent\textbf{Q47. A designer connects a 3.3 V logic sensor's output to an Arduino Uno ( 5 V logic) input and reads correct HIGH values, but connecting the Uno's 5 V output to the sensor's input risks damage. What is the correct explanation?} \hfill [Topic: introduction of arduino and sensors | Hard]
\begin{enumerate}[label=(\alph*),topsep=2pt,itemsep=1pt,partopsep=0pt,parsep=0pt]
    \item {Both directions are always safe without level shifting}
    \item {3.3 V exceeds the Uno's HIGH threshold, but 5 V exceeds the sensor's max input rating}
    \item {The Uno interprets 3.3 V as LOW always}
    \item {The sensor cannot output any logic HIGH at 3.3 V}
\end{enumerate}
\vspace{2mm}

\noindent\textbf{Q48. A resistive sensor is read via a voltage divider with a fixed 10 k $\Omega$ resistor and 5 V supply. For best ADC sensitivity near the sensor's nominal resistance R s = 10 k $\Omega$ , which statement is correct?} \hfill [Topic: introduction of arduino and sensors | Hard]
\begin{enumerate}[label=(\alph*),topsep=2pt,itemsep=1pt,partopsep=0pt,parsep=0pt]
    \item {Output voltage is independent of R s}
    \item {The ADC cannot resolve any change at this point}
    \item {Sensitivity is maximized because the divider output is near mid-scale ( 2.5 V )}
    \item {Sensitivity is lowest because output saturates at 5 V}
\end{enumerate}
\vspace{2mm}

\noindent\textbf{Q49. Why does the Arduino delay() function create problems in a system that must simultaneously poll three sensors at different rates?} \hfill [Topic: introduction of arduino and sensors | Hard]
\begin{enumerate}[label=(\alph*),topsep=2pt,itemsep=1pt,partopsep=0pt,parsep=0pt]
    \item {It increases the ADC resolution unpredictably}
    \item {It resets the microcontroller after each call}
    \item {It disables all GPIO pins permanently}
    \item {It blocks the CPU, preventing other tasks during the delay interval}
\end{enumerate}
\vspace{2mm}

\noindent\textbf{Q50. In a plant condition-monitoring system, vibration sensors sample at 10 kHz . To detect a fault signature at 4 kHz without aliasing, which sampling and processing choice is correct?} \hfill [Topic: role of IOT and cloud computing in condition monitoring of plant processes | Hard]
\begin{enumerate}[label=(\alph*),topsep=2pt,itemsep=1pt,partopsep=0pt,parsep=0pt]
    \item {10 kHz is below Nyquist for 4 kHz , so detection fails}
    \item {Sampling frequency is irrelevant to aliasing}
    \item {A 4 kHz signal requires at least 4 kHz sampling only}
    \item {10 kHz satisfies Nyquist ( > 2 $\times$ 4 kHz ), so 4 kHz is detectable}
\end{enumerate}
\vspace{2mm}

\noindent\textbf{Q51. A plant uses edge computing before sending data to the cloud. Which scenario BEST justifies processing at the edge rather than purely in the cloud?} \hfill [Topic: role of IOT and cloud computing in condition monitoring of plant processes | Hard]
\begin{enumerate}[label=(\alph*),topsep=2pt,itemsep=1pt,partopsep=0pt,parsep=0pt]
    \item {Cross-site statistical benchmarking of all plants}
    \item {Monthly aggregate report generation}
    \item {Long-term trend archival over several years}
    \item {Low-latency shutdown decisions needed within milliseconds on local vibration spikes}
\end{enumerate}
\vspace{2mm}

\noindent\textbf{Q52. An IoT gateway transmits 500 sensor readings per second, each 20 bytes , over a link. What minimum raw throughput is required, ignoring protocol overhead?} \hfill [Topic: role of IOT and cloud computing in condition monitoring of plant processes | Hard]
\begin{enumerate}[label=(\alph*),topsep=2pt,itemsep=1pt,partopsep=0pt,parsep=0pt]
    \item {40 kbps}
    \item {80 kbps}
    \item {160 kbps}
    \item {10 kbps}
\end{enumerate}
\vspace{2mm}

\noindent\textbf{Q53. A wearable IoT ECG device must prioritize detecting life-threatening arrhythmia locally while uploading full waveforms to the cloud. Which architecture best balances reliability and bandwidth?} \hfill [Topic: health care | Hard]
\begin{enumerate}[label=(\alph*),topsep=2pt,itemsep=1pt,partopsep=0pt,parsep=0pt]
    \item {Run arrhythmia detection on-device; stream compressed waveforms to cloud for review}
    \item {Send every raw sample to cloud and wait for cloud diagnosis before alerting}
    \item {Store all data on-device only and never transmit}
    \item {Perform all detection in the cloud with no local processing}
\end{enumerate}
\vspace{2mm}

\noindent\textbf{Q54. In a cloud-connected patient monitoring system, which design flaw most directly violates patient data security requirements?} \hfill [Topic: health care | Hard]
\begin{enumerate}[label=(\alph*),topsep=2pt,itemsep=1pt,partopsep=0pt,parsep=0pt]
    \item {Using redundant cloud storage regions}
    \item {Buffering readings locally before upload}
    \item {Transmitting physiological data over unencrypted HTTP}
    \item {Compressing waveforms before transmission}
\end{enumerate}
\vspace{2mm}

\noindent\textbf{Q55. A smart irrigation system uses soil-moisture sensors and a cloud model. During a multi-day network outage, which edge-level fallback best preserves crop health?} \hfill [Topic: agriculture | Hard]
\begin{enumerate}[label=(\alph*),topsep=2pt,itemsep=1pt,partopsep=0pt,parsep=0pt]
    \item {System irrigates continuously at maximum flow}
    \item {Local controller irrigates using threshold rules on cached sensor data}
    \item {System deletes sensor data to save memory}
    \item {System halts all irrigation until cloud reconnects}
\end{enumerate}
\vspace{2mm}

\noindent\textbf{Q56. A field deploys 200 LoRa soil sensors each reporting once every 10 minutes . What is the average aggregate message rate the gateway must handle?} \hfill [Topic: agriculture | Hard]
\begin{enumerate}[label=(\alph*),topsep=2pt,itemsep=1pt,partopsep=0pt,parsep=0pt]
    \item {$\approx$ 0.33 messages/s}
    \item {$\approx$ 20 messages/s}
    \item {$\approx$ 3.3 messages/s}
    \item {$\approx$ 2 messages/s}
\end{enumerate}
\vspace{2mm}

\noindent\textbf{Q57. In a smart factory, predictive maintenance uses IoT vibration data and cloud ML. What is the primary advantage over traditional scheduled (time-based) maintenance?} \hfill [Topic: manufacturing | Hard]
\begin{enumerate}[label=(\alph*),topsep=2pt,itemsep=1pt,partopsep=0pt,parsep=0pt]
    \item {Maintenance is triggered by actual equipment condition, reducing unplanned downtime}
    \item {It removes the need for cloud or network connectivity}
    \item {It guarantees machines never fail under any circumstance}
    \item {It eliminates the need for any sensors on machines}
\end{enumerate}
\vspace{2mm}

\noindent\textbf{Q58. A manufacturing line's IoT system shows rising false-alarm rates after adding more sensors. Which root cause is most consistent with this symptom?} \hfill [Topic: manufacturing | Hard]
\begin{enumerate}[label=(\alph*),topsep=2pt,itemsep=1pt,partopsep=0pt,parsep=0pt]
    \item {Removal of the cloud analytics layer}
    \item {Reduced total data volume from the line}
    \item {Fixed thresholds not adapted to increased sensor noise and correlation}
    \item {Lower sampling frequency on all channels}
\end{enumerate}
\vspace{2mm}

\noindent\textbf{Q59. In a connected vehicle, which function MUST remain on an embedded ECU rather than depend on cloud computing, and why?} \hfill [Topic: automobiles | Hard]
\begin{enumerate}[label=(\alph*),topsep=2pt,itemsep=1pt,partopsep=0pt,parsep=0pt]
    \item {Over-the-air firmware distribution, due to version management}
    \item {Infotainment music streaming, due to large data size}
    \item {Anti-lock braking control, due to millisecond-level safety-critical latency}
    \item {Fleet-wide fuel-economy benchmarking, due to data aggregation}
\end{enumerate}
\vspace{2mm}

\noindent\textbf{Q60. A smart grid uses IoT smart meters for demand response. During a sudden generation shortfall, which cloud+edge mechanism most effectively prevents a cascading blackout?} \hfill [Topic: smart grid | Hard]
\begin{enumerate}[label=(\alph*),topsep=2pt,itemsep=1pt,partopsep=0pt,parsep=0pt]
    \item {Disabling all smart meters until supply recovers}
    \item {Increasing data logging resolution only}
    \item {Edge-level automatic load shedding guided by cloud demand-response signals}
    \item {Waiting for monthly cloud billing reports to adjust tariffs}
\end{enumerate}
\vspace{2mm}


\newpage
\section{Section B: Comprehensive Analytical Problems with Step-by-Step Solutions}

\begin{subjectiveq}{6}{1}
An Arduino Uno microcontroller development system based on the ATmega328P is deployed to monitor agricultural conditions using an analog temperature sensor (LM35) and a soil moisture sensor.\begin{enumerate}[label=(\alph*)]\item Contrast the architecture of the ATmega328P microcontroller against an Intel Core i7 microprocessor across memory organization, power consumption, execution determinism, and peripheral integration.\item The ATmega328P features a 10-bit Successive Approximation ADC operating with reference voltage $V_{ref} = 5.0\text{ V}$. Compute the quantization step size. If the LM35 produces $10\text{ mV}/^\circ\text{C}$ and reads $35.0^\circ\text{C}$, determine the analog voltage and the digital ADC code recorded.\item Outline the end-to-end cloud IoT architecture for telemetry streaming from this sensor node to a remote cloud analytics dashboard via MQTT/Wi-Fi.\end{enumerate}
\end{subjectiveq}

\begin{solbox}
\textbf{Step 1: Architectural Comparison (Microcontroller vs Microprocessor)}\\\begin{center}\small\begin{tabularx}{\linewidth}{l X X}\toprule\textbf{Feature} & \textbf{ATmega328P (Microcontroller)} & \textbf{Intel Core i7 (Microprocessor)} \\\midruleArchitecture & Harvard (Separate code and data buses) & Modified Harvard / Unified cache \\Hardware On-Chip & CPU, 32KB Flash, 2KB SRAM, ADC, Timers, UART on single chip & CPU cores, cache only; requires external RAM, GPU, Chipset \\Power & Very low ($< 50\text{ mW}$) & High (65--125 W) \\Determinism & Hard real-time cycle-accurate execution & Non-deterministic OS latency \\Typical Domain & Dedicated sensing and actuator control & General-purpose computing systems \\\bottomrule\end{tabularx}\end{center}\vspace{2mm}\textbf{Step 2: ADC Quantization and Sensor Signal Digital Conversion}\\A 10-bit ADC provides $2^{10} = 1024$ discrete digital levels ($0$ to $1023$).The quantization step size ($V_{LSB}$) is:\[ V_{LSB} = \frac{V_{ref}}{1024} = \frac{5.0\text{ V}}{1024} \approx 4.8828\text{ mV} \]At $T = 35.0^\circ\text{C}$, the LM35 output voltage is $V_{out} = 35.0 \times 10\text{ mV} = 350.0\text{ mV} = 0.350\text{ V}$.The digital ADC count is:\[ D = \left\lfloor \frac{0.350}{5.0} \times 1024 \right\rfloor = \lfloor 71.68 \rfloor = 71 \text{ or } 72 \]\textbf{Step 3: End-to-End Cloud IoT Telemetry Architecture}\\\begin{enumerate}[leftmargin=1.5em]\item \textbf{Perception Layer}: Sensors convert physical quantities into analog voltages.\item \textbf{Edge Layer}: ATmega328P digitizes and serializes measurements into JSON packets.\item \textbf{Network Layer}: Wi-Fi module transmits packets using MQTT over TLS to the cloud broker.\item \textbf{Application Layer}: Cloud service ingests, stores in time-series DB, and displays real-time telemetry.\end{enumerate}
\end{solbox}

\begin{examinertip}
\textbf{Implementation Note}: Using the internal $1.1\text{ V}$ reference increases LM35 sensitivity to $\approx 1.074\text{ mV/step}$.
\end{examinertip}
\vspace{4mm}
